% ============================================================================
% Refmon: pipeline shapes, diagrams and flowcharts, level by level.
% Build:  make diagrams      (or: pdflatex diagrams.tex, twice)
% Every figure is a plain tikzpicture in diagrams/<name>.tex, so main.tex can
% reuse any of them:   % ============================================================================
% Refmon diagram styles: shared by every file in documentation/diagrams/.
% Needs \usepackage{tikz} and xcolor. It is safe to \input this after
% main.tex's own \usetikzlibrary line: loading a library twice is harmless,
% and every style and colour here carries an `rm' prefix so nothing in
% main.tex is redefined.
% ============================================================================
\usetikzlibrary{arrows.meta,positioning,shapes.geometric,shapes.multipart,%
  shapes.symbols,calc,fit,backgrounds,decorations.pathreplacing}

% ---- palette (light, print-safe) --------------------------------------------
\definecolor{rmInk}{HTML}{24292F}
\definecolor{rmStage}{HTML}{DDF4FF}      \definecolor{rmStageLine}{HTML}{0969DA}
\definecolor{rmStore}{HTML}{EAEEF2}      \definecolor{rmStoreLine}{HTML}{57606A}
\definecolor{rmDecide}{HTML}{FFF8C5}     \definecolor{rmDecideLine}{HTML}{9A6700}
\definecolor{rmAllow}{HTML}{DAFBE1}      \definecolor{rmAllowLine}{HTML}{1A7F37}
\definecolor{rmHold}{HTML}{FFF1E5}       \definecolor{rmHoldLine}{HTML}{BC4C00}
\definecolor{rmBlock}{HTML}{FFEBE9}      \definecolor{rmBlockLine}{HTML}{CF222E}
\definecolor{rmHuman}{HTML}{FBEFFF}      \definecolor{rmHumanLine}{HTML}{8250DF}

% ---- signal chips used inside node text ------------------------------------
% \rmC{contract\_x}  \rmE{effect\_x}  \rmF{flag}   (escape underscores!)
\providecommand{\rmChip}[2]{{\setlength{\fboxsep}{0.8pt}\colorbox{#1}{\strut\texttt{#2}}}}
\providecommand{\rmC}[1]{\rmChip{rmStage}{#1}}
\providecommand{\rmE}[1]{\rmChip{rmHold}{#1}}
\providecommand{\rmF}[1]{\rmChip{rmBlock}{#1}}
% U(allowance, effect): the universal-contract helper of signals.py
\providecommand{\rmU}[2]{U(\texttt{#1},\,\rmE{#2})}

\tikzset{
  % picture-level defaults: \begin{tikzpicture}[rm]
  rm/.style={font=\scriptsize\sffamily, align=center, line width=0.55pt,
             >={Latex[length=1.6mm,width=1.3mm]}},
  % nodes
  rm stage/.style={rectangle, rounded corners=2.5pt, draw=rmStageLine, fill=rmStage,
             thick, minimum height=9mm, minimum width=12mm, inner sep=2pt},
  rm box/.style={rectangle, draw=rmInk!60, fill=white, minimum height=6mm,
             minimum width=14mm, inner sep=2.2pt},
  rm sub/.style={rm box, draw=rmStageLine!70, fill=rmStage!40},
  rm dec/.style={diamond, aspect=2.1, draw=rmDecideLine, fill=rmDecide, thick,
             inner sep=0.5pt, font=\tiny\sffamily},
  rm term/.style={rectangle, rounded corners=2.6mm, minimum height=5.4mm,
             minimum width=14mm, draw=rmInk!60, fill=rmStore, inner sep=2pt},
  rm store/.style={cylinder, shape border rotate=90, aspect=0.2, draw=rmStoreLine,
             fill=rmStore, minimum height=9mm, minimum width=15mm, inner sep=1.5pt},
  rm io/.style={trapezium, trapezium left angle=75, trapezium right angle=105,
             trapezium stretches=true, draw=rmInk!60, fill=white,
             minimum height=7mm, minimum width=16mm, inner sep=1.5pt},
  rm human/.style={rectangle, rounded corners=4pt, draw=rmHumanLine, fill=rmHuman,
             thick, minimum height=7mm, minimum width=16mm, inner sep=2pt},
  rm raw/.style={rm box, draw=rmBlockLine, fill=rmBlock},
  rm allow/.style={rm term, draw=rmAllowLine, fill=rmAllow, thick,
             font=\scriptsize\sffamily\bfseries},
  rm hold/.style={rm term, draw=rmHoldLine, fill=rmHold, thick,
             font=\scriptsize\sffamily\bfseries},
  rm block/.style={rm term, draw=rmBlockLine, fill=rmBlock, thick,
             font=\scriptsize\sffamily\bfseries},
  rm chev/.style={signal, signal to=east, signal from=west, draw=rmStageLine,
             fill=rmStage, minimum height=9mm, minimum width=17mm, inner sep=1pt,
             font=\tiny\sffamily},
  % edges
  rm flow/.style={->, draw=rmInk!85},
  rm back/.style={->, dashed, draw=rmStageLine, thick},
  rm ctrl/.style={->, densely dotted, draw=rmStoreLine, thick},
  rm hflow/.style={->, draw=rmHumanLine, thick},
  rm bad/.style={->, draw=rmBlockLine, thick},
  % labels
  rm yes/.style={font=\tiny\sffamily\bfseries, text=rmAllowLine, inner sep=1pt},
  rm no/.style={font=\tiny\sffamily\bfseries, text=rmBlockLine, inner sep=1pt},
  rm note/.style={font=\tiny\sffamily, text=rmInk!80, inner sep=1pt},
  rm lbl/.style={font=\tiny\sffamily, fill=white, inner sep=1pt},
  % groups (draw on the background layer with fit=...)
  rm group/.style={draw=rmStoreLine!70, dashed, rounded corners=5pt, inner sep=2.5mm},
  rm lane/.style={draw=none, fill=#1, rounded corners=4pt, inner sep=2.5mm},
  rm group label/.style={font=\tiny\sffamily\bfseries, text=rmStoreLine,
             inner sep=1.5pt, anchor=south west},
  rm wall/.style={draw=rmBlockLine, very thick, dashed},
}
  (once, in the preamble)
%                      \begin{figure*}[!t]\centering
%                        \adjustbox{max width=\textwidth}{% Level 1: one call of Pipeline.process(event) (src/refmon/pipeline/pipeline.py). Wide figure.
\begin{tikzpicture}[rm]
  % ---- row 1: the stages -------------------------------------------------
  \node[rm io, minimum width=24mm] (in) at (-0.4,0) {\texttt{RawEvent} | \texttt{Event}\\\tiny from \texttt{Sensor.emit}};
  \node[rm stage, minimum width=17mm] (obs) at (2.15,0)  {\textbf{Observe}\\\tiny fingerprint, drop raw};
  \node[rm stage, minimum width=17mm] (tag) at (4.25,0)  {\textbf{Tag}\\\tiny rules 1--7};
  \node[rm stage, minimum width=17mm] (sig) at (6.35,0)  {\textbf{Signals}\\\tiny \texttt{extract()}};
  \node[rm stage, minimum width=17mm] (lnk) at (8.45,0)  {\textbf{Link}\\\tiny assign / merge story};
  \node[rm stage, minimum width=17mm] (jdg) at (10.55,0) {\textbf{Judge}\\\tiny $\to$ \texttt{Judgement}};
  \draw[rm flow] (in) -- (obs);
  \draw[rm flow] (obs) -- (tag);
  \draw[rm flow] (tag) -- (sig);
  \draw[rm flow] (sig) -- (lnk);
  \draw[rm flow] (lnk) -- (jdg);
  \draw[rm wall] (3.2,-0.65) -- (3.2,0.65);
  \node[rm note, text=rmBlockLine, below] at (3.2,-0.65) {privacy wall};

  % ---- state the stages share ---------------------------------------------
  \node[rm store] (buf) at (2.15,1.75) {Ring buffer\\\tiny Events only};
  \node[rm store] (tst) at (4.25,1.75) {Tagger state\\\tiny labels, procs};
  \node[rm store] (rts) at (6.35,1.75) {Rates\\\tiny windows, EMA};
  \draw[rm ctrl] (obs) -- (buf);
  \draw[rm ctrl, <->] (tag) -- (tst);
  \draw[rm ctrl] (rts) -- (sig);

  % ---- row 2: after the judgement ----------------------------------------
  \node[rm box] (inc) at (10.9,-1.9) {\texttt{rates.mark\_incident}\\\tiny freeze baselines while $D\neq$ allow};
  \node[rm box] (evd) at (7.6,-1.9)  {\texttt{tagger.set\_evidence}\\\tiny labels behind a Contract\\\tiny signal stop decaying (D21, A49)};
  \node[rm dec] (emd) at (4.75,-1.9) {decision\\changed and\\$\neq$ ALLOW?};
  \node[rm box] (vrd) at (1.75,-1.9) {emit \texttt{Verdict}\\\tiny contributions, $L$, threshold, gate};
  \draw[rm flow] (jdg.south) -- (jdg.south |- inc.north);
  \draw[rm flow] (inc) -- (evd);
  \draw[rm flow] (evd) -- (emd);
  \draw[rm flow] (emd) -- node[rm yes, above] {yes} (vrd);

  % feedback loops into shared state
  \draw[rm back] (evd.north) |- (5.3,-1.0) |- (tst.east);
  \node[rm note, text=rmStageLine, anchor=west] at (5.35,-0.75) {labels frozen};
  \draw[rm back] (inc.east) -- ++(0.35,0) |- (rts.east);
  \node[rm note, text=rmStageLine, anchor=south east] at (12.5,1.75) {baselines frozen};

  % ---- row 3: the Gate ----------------------------------------------------
  \node[rm dec] (gq)   at (4.75,-3.75) {Gate\\enabled?};
  \node[rm box] (tick) at (7.6,-3.75)  {\texttt{Gate.tick(ts)}\\\tiny resolve expired holds};
  \node[rm box] (hnd)  at (10.9,-3.75) {\texttt{Gate.handle}\\\tiny hold / block via \texttt{Contain.plan}};
  \node[rm term] (ret) at (4.75,-5.15) {return \texttt{[*tick, *verdicts, *actions]}};
  \draw[rm flow] (emd) -- node[rm no, right] {no} (gq);
  \draw[rm flow] (vrd.south) |- (gq.west);
  \draw[rm flow] (gq) -- node[rm yes, above] {yes} (tick);
  \draw[rm flow] (tick) -- (hnd);
  \draw[rm flow] (gq) -- node[rm no, right] {no} (ret);
  \draw[rm flow] (hnd.south) |- (ret.east);
\end{tikzpicture}
}
%                        \caption{...}\end{figure*}
% ============================================================================
\documentclass[10pt,letterpaper]{article}
\usepackage[utf8]{inputenc}
\usepackage[T1]{fontenc}
\usepackage{lmodern}
\usepackage{microtype}
\usepackage[margin=0.72in]{geometry}
\usepackage{amsmath,amssymb}
\usepackage{xcolor}
\usepackage{tikz}
\usepackage{adjustbox}
\usepackage[colorlinks=true,linkcolor=blue!60!black,urlcolor=blue!65!black]{hyperref}

% ============================================================================
% Refmon diagram styles: shared by every file in documentation/diagrams/.
% Needs \usepackage{tikz} and xcolor. It is safe to \input this after
% main.tex's own \usetikzlibrary line: loading a library twice is harmless,
% and every style and colour here carries an `rm' prefix so nothing in
% main.tex is redefined.
% ============================================================================
\usetikzlibrary{arrows.meta,positioning,shapes.geometric,shapes.multipart,%
  shapes.symbols,calc,fit,backgrounds,decorations.pathreplacing}

% ---- palette (light, print-safe) --------------------------------------------
\definecolor{rmInk}{HTML}{24292F}
\definecolor{rmStage}{HTML}{DDF4FF}      \definecolor{rmStageLine}{HTML}{0969DA}
\definecolor{rmStore}{HTML}{EAEEF2}      \definecolor{rmStoreLine}{HTML}{57606A}
\definecolor{rmDecide}{HTML}{FFF8C5}     \definecolor{rmDecideLine}{HTML}{9A6700}
\definecolor{rmAllow}{HTML}{DAFBE1}      \definecolor{rmAllowLine}{HTML}{1A7F37}
\definecolor{rmHold}{HTML}{FFF1E5}       \definecolor{rmHoldLine}{HTML}{BC4C00}
\definecolor{rmBlock}{HTML}{FFEBE9}      \definecolor{rmBlockLine}{HTML}{CF222E}
\definecolor{rmHuman}{HTML}{FBEFFF}      \definecolor{rmHumanLine}{HTML}{8250DF}

% ---- signal chips used inside node text ------------------------------------
% \rmC{contract\_x}  \rmE{effect\_x}  \rmF{flag}   (escape underscores!)
\providecommand{\rmChip}[2]{{\setlength{\fboxsep}{0.8pt}\colorbox{#1}{\strut\texttt{#2}}}}
\providecommand{\rmC}[1]{\rmChip{rmStage}{#1}}
\providecommand{\rmE}[1]{\rmChip{rmHold}{#1}}
\providecommand{\rmF}[1]{\rmChip{rmBlock}{#1}}
% U(allowance, effect): the universal-contract helper of signals.py
\providecommand{\rmU}[2]{U(\texttt{#1},\,\rmE{#2})}

\tikzset{
  % picture-level defaults: \begin{tikzpicture}[rm]
  rm/.style={font=\scriptsize\sffamily, align=center, line width=0.55pt,
             >={Latex[length=1.6mm,width=1.3mm]}},
  % nodes
  rm stage/.style={rectangle, rounded corners=2.5pt, draw=rmStageLine, fill=rmStage,
             thick, minimum height=9mm, minimum width=12mm, inner sep=2pt},
  rm box/.style={rectangle, draw=rmInk!60, fill=white, minimum height=6mm,
             minimum width=14mm, inner sep=2.2pt},
  rm sub/.style={rm box, draw=rmStageLine!70, fill=rmStage!40},
  rm dec/.style={diamond, aspect=2.1, draw=rmDecideLine, fill=rmDecide, thick,
             inner sep=0.5pt, font=\tiny\sffamily},
  rm term/.style={rectangle, rounded corners=2.6mm, minimum height=5.4mm,
             minimum width=14mm, draw=rmInk!60, fill=rmStore, inner sep=2pt},
  rm store/.style={cylinder, shape border rotate=90, aspect=0.2, draw=rmStoreLine,
             fill=rmStore, minimum height=9mm, minimum width=15mm, inner sep=1.5pt},
  rm io/.style={trapezium, trapezium left angle=75, trapezium right angle=105,
             trapezium stretches=true, draw=rmInk!60, fill=white,
             minimum height=7mm, minimum width=16mm, inner sep=1.5pt},
  rm human/.style={rectangle, rounded corners=4pt, draw=rmHumanLine, fill=rmHuman,
             thick, minimum height=7mm, minimum width=16mm, inner sep=2pt},
  rm raw/.style={rm box, draw=rmBlockLine, fill=rmBlock},
  rm allow/.style={rm term, draw=rmAllowLine, fill=rmAllow, thick,
             font=\scriptsize\sffamily\bfseries},
  rm hold/.style={rm term, draw=rmHoldLine, fill=rmHold, thick,
             font=\scriptsize\sffamily\bfseries},
  rm block/.style={rm term, draw=rmBlockLine, fill=rmBlock, thick,
             font=\scriptsize\sffamily\bfseries},
  rm chev/.style={signal, signal to=east, signal from=west, draw=rmStageLine,
             fill=rmStage, minimum height=9mm, minimum width=17mm, inner sep=1pt,
             font=\tiny\sffamily},
  % edges
  rm flow/.style={->, draw=rmInk!85},
  rm back/.style={->, dashed, draw=rmStageLine, thick},
  rm ctrl/.style={->, densely dotted, draw=rmStoreLine, thick},
  rm hflow/.style={->, draw=rmHumanLine, thick},
  rm bad/.style={->, draw=rmBlockLine, thick},
  % labels
  rm yes/.style={font=\tiny\sffamily\bfseries, text=rmAllowLine, inner sep=1pt},
  rm no/.style={font=\tiny\sffamily\bfseries, text=rmBlockLine, inner sep=1pt},
  rm note/.style={font=\tiny\sffamily, text=rmInk!80, inner sep=1pt},
  rm lbl/.style={font=\tiny\sffamily, fill=white, inner sep=1pt},
  % groups (draw on the background layer with fit=...)
  rm group/.style={draw=rmStoreLine!70, dashed, rounded corners=5pt, inner sep=2.5mm},
  rm lane/.style={draw=none, fill=#1, rounded corners=4pt, inner sep=2.5mm},
  rm group label/.style={font=\tiny\sffamily\bfseries, text=rmStoreLine,
             inner sep=1.5pt, anchor=south west},
  rm wall/.style={draw=rmBlockLine, very thick, dashed},
}


% \rmdiagram{file}{caption}  ->  figure with label fig:<file>, scaled down only if too large
\newcommand{\rmdiagram}[2]{%
  \begin{figure}[!htbp]
    \centering
    \adjustbox{max width=\linewidth, max totalheight=0.86\textheight}{\input{diagrams/#1}}%
    \caption{#2}\label{fig:#1}
  \end{figure}}

\setlength{\parskip}{3pt}
\setlength{\parindent}{0pt}

\title{\bfseries Refmon (Provenance Gate v2)\\[2pt]\large Pipeline shapes, diagrams and flowcharts, level by level}
\author{Companion to the Refmon paper (\texttt{main.tex}); drawn from the source in \texttt{src/refmon/}}
\date{Specification tag \texttt{v1.3}}

\begin{document}
\maketitle

\noindent The figures go from the outside in. \textbf{Level~0} places Refmon among sensors, enforcers and
people; \textbf{Level~1} follows one event through \texttt{Pipeline.process}; \textbf{Level~2} opens every
stage; \textbf{Level~3} shows how the policy is governed and what leaves a device; \textbf{Level~4} shows
how the results are produced; \textbf{Level~5} shows the live user-space deployment. Every node names the
module or function it stands for, and every figure is a stand-alone \texttt{tikzpicture} in
\texttt{diagrams/}.

\rmdiagram{legend}{Shapes and arrows used in every figure. Blue chips are Contract signals, orange chips
Effect signals, red chips override flags.}

\tableofcontents
\listoffigures
\clearpage

% ============================================================================
\section{Level 0: system context}
Only two interfaces touch a platform: \texttt{Sensor.emit(sink)} and \texttt{Enforcer.act(Action)}
(CLAUDE.md C1). Everything between them is platform-independent, and raw content never crosses the
privacy wall after \texttt{Observe}.

\rmdiagram{L0-system-context}{System context. Five kinds of sensor feed the platform-independent core;
the signed policy, the Exposure overlay and the self-test sit beside it on the device. Actions go to the
enforcer adapter, hold prompts to people, and only pseudonymized verdicts, heartbeats and noisy counts
leave the device.}

% ============================================================================
\clearpage
\section{Level 1: one event through the core}
\rmdiagram{L1-core-pipeline}{One call of \texttt{Pipeline.process(event)}. After the judgement, the
rate baselines are frozen while the story is not allowed, and the labels behind a Contract signal stop
decaying (dashed feedback into shared state). A verdict is emitted only when the decision changed and is
not ALLOW; the Gate then resolves expired holds and acts on the story.}

% ============================================================================
\clearpage
\section{Level 2: inside each stage}
\subsection{Observe}
\rmdiagram{L2a-observe}{Observe, the only stage that sees raw content. Outcome attributes a sensor sets
itself are discarded; content becomes a keyed HMAC fingerprint; obfuscation features are kept as
numbers; every decision that needs a raw host, URL, program or signer is made here; then the raw values
are dropped.}

\subsection{Tag}
\rmdiagram{L2b-tag-rules}{Tag: propagation rules 1--7 as a dispatch on the event type. A ``no'' at any
other decision leaves the labels unchanged. Rule 7 (persistent labels carry Contract evidence) is applied
by the pipeline after the judgement (Figure~\ref{fig:L1-core-pipeline}).}

\rmdiagram{L2b2-label-strength}{(a) Label strength over time: an ordinary label halves every 30 minutes
and stops failing Origin at $\theta_{ori}=0.5$; a persistent label never decays; a label behind a
Contract signal is frozen at the strength it had (A49). (b) The F8 fix: up to v1.2 a Contract or Effect
signal counted at the label's strength above $\theta_{ori}$ and at full strength below it, so a weaker
label could give more evidence; since v1.3 it always counts at full strength.}

\rmdiagram{L2c-provenance-S06}{A worked provenance graph, golden scenario S06 (persistence now, payload a
day later). The startup item gets a persistent label L2 that records the Contract signals of the first
story (D21), so the delayed run, a new story with no paste in it, still fails all three questions.}

\subsection{Link}
\rmdiagram{L2d-link}{Link: an event joins every open story that shares a key with it within the quiet
window; several such stories are merged.}

\subsection{Signals}
\rmdiagram{L2e-signals}{Signal extraction per event type (\texttt{signals.py}). The prefix of a signal
name is its question; the weight of each signal in nats is in the signed policy.}

\subsection{Judge and the evidence calculus}
\rmdiagram{L2f-judge}{Judge: effective signals, overrides in priority order, the plain 3/2/1 rule, and
the evidence calculus with the contract gate.}

\rmdiagram{L2g-calculus-waterfall}{The calculus with the default weights. Each column starts at the prior
$L_0=-6.91$ nats and adds the largest contribution of each question; the bands are the cost-derived
thresholds. The last column is illustrative: it crosses $\tau_b$ without a Contract signal, so the
contract gate caps it at HOLD.}

\rmdiagram{L2h-rates}{Rate tests behind \texttt{mass\_change}, \texttt{lateral}, repeated
\texttt{response\_anomaly} (Poisson tail over a short and a long window) and \texttt{exfil} (CUSUM over
standardized volume). All baselines freeze while an incident is open.}

\subsection{Gate and Contain}
\rmdiagram{L2i-gate}{\texttt{Gate.handle} for one judged story: the mode downgrade (observe never acts,
warn never blocks), blocks, and the start or extension of a hold. Expired holds follow the hold state
machine of the paper (endpoint: stop; server: stay quarantined).}

\rmdiagram{L2j-contain}{\texttt{Contain.plan}: only labeled processes and what they wrote are touched;
when nothing is labeled, the processes that produced the signals are, but never a channel role.}

% ============================================================================
\clearpage
\section{Level 3: governance and privacy}
\rmdiagram{L3a-policy-lifecycle}{The signed policy from key to running pipeline. Any failed check keeps the
last valid policy; the Exposure overlay only tightens it.}

\rmdiagram{L3b-learning-tune}{How the policy may change: learning-mode contract proposals and service-desk
feedback become proposals, a named person approves chosen changes, and the result still has to be signed.}

\rmdiagram{L3c-exposure}{Exposure stage: every component gets one state, and an end-of-life or exposed
component escalates its own role only (virtual patching). Scenarios S16, S28 and S29.}

\rmdiagram{L3d-selftest-heartbeat}{Continuous self-test in an isolated copy of the live pipeline, and the
heartbeat that reports switched-off protections and degraded sensor capabilities.}

\rmdiagram{L3e-privacy-reporting}{What leaves a device and how: plain-language reports under pseudonyms,
dashboard counts with $k$-anonymity and Laplace noise, re-identification only with two approvers and an
audit log, and the adversarial privacy audit over all outputs.}

% ============================================================================
\clearpage
\section{Level 4: evaluation}
\rmdiagram{L4a-reproduce}{\texttt{make reproduce}: one seeded variant file replayed in every tagged
version, compared, checked against the committed per-variant data, and written into
\texttt{docs/results.md} together with the golden, coverage, audit, publishability and test results.}

\rmdiagram{L4b-phases-versions}{Build phases (top) and the version history (bottom): each version answers
the findings of the previous evaluation and is measured on a seed recorded before the change.
``Contained'' is block plus effective hold on the 633 malicious variants of seed 654514256.}

\rmdiagram{L4c-test-layers}{The six test layers of the specification (C9).}

% ============================================================================
\clearpage
\section{Level 5: live user-space deployment}
\rmdiagram{L5a-live-guard}{The live command barrier (\texttt{refmon guard}). The command is turned into
synthetic raw events, judged by the unchanged pipeline, and run only if the barrier allows it.}

\rmdiagram{L5b-sandbox-watcher}{The sandbox watcher of \texttt{refmon serve}: files dropped into the web
root are judged as web-worker writes and quarantined on a block or hold; workspace bursts are judged as
mass change.}

\rmdiagram{L5c-rootkit}{The read-only rootkit scanner. Cross-view engines are intersected over $m$ passes
so transient churn drops out; findings are fingerprinted, then scored twice: directly by the calculus and
through the pipeline.}

\end{document}
